\documentclass[11pt]{article}
\usepackage{fullpage}
%% color for the links 
\usepackage[usenames,dvipsnames]{color}
\usepackage{longtable, graphicx}
\usepackage{multicol}
\usepackage{ebgaramond}		% for boldface small caps
\usepackage{changepage}		% for indented paragraph blocks
% We will generate all images so they have a width \maxwidth. This means
% that they will get their normal width if they fit onto the page, but
% are scaled down if they would overflow the margins.
\makeatletter
\def\maxwidth{\ifdim\Gin@nat@width>\linewidth\linewidth
\else\Gin@nat@width\fi}
\makeatother
\let\Oldincludegraphics\includegraphics
% \renewcommand{\includegraphics}[1]{\Oldincludegraphics[width=\maxwidth]{#1}}
%% hyperlinks
\usepackage[
	colorlinks=true,
	urlcolor=MidnightBlue,
	plainpages=false,
  	]{hyperref}% color for the links 
\setlength{\parindent}{0pt}
\setlength{\parskip}{6pt plus 2pt minus 1pt}
\setlength{\emergencystretch}{3em}  % prevent overfull lines
\setcounter{secnumdepth}{0}

\begin{document}

\begin{multicols}{2}

	\includegraphics[width = 0.45\textwidth]{../../images/rice_smgt.png}

	\columnbreak

	\raggedleft
	{\bf Course:} SMGT $435$ -- Baseball Analytics\\
	{\bf Term:} Spring $2024$\\
	{\bf Class:} $2$:$30$ p.m. -- $3$:$45$ p.m. Tue/Thu

\end{multicols}

\hrule

~\\
\textbf{\textsc{Instructor Contact Information}}
\begin{adjustwidth}{1cm}{0cm}
	{\bf Instructor:} Scott Powers\\
  {\bf Office:} Kraft $246$\\
  {\bf Email:} scott.powers@rice.edu\\
  {\bf Office Hours:} https://calendly.com/saberpowers
\end{adjustwidth}

~\\
\textbf{\textsc{Required Texts and Materials}}
\begin{adjustwidth}{1cm}{0cm}
	{\bf Class Website:} https://canvas.rice.edu\\
	{\bf Required Text:} All readings and videos will be made available on Canvas.\\
	{\bf Required Material:} Bring pen/pencil and paper to class every day. Bring your laptop on ThuRsdays.
\end{adjustwidth}

~\\
\textbf{\textsc{Topics}}
\begin{adjustwidth}{1cm}{0cm}
  Pythagorean Formula, Base-Out Run Expectancy, Linear Weights, Batted Ball Outcome Model, Fielding, Baserunning, BABIP/FIP/DIPS, Count Value, Catcher Framing, Pitch Outcome Model, Replacement Level, Positional Adjustment, WAR, Projections, Aging Curves, Level Translations, Additional Topics
\end{adjustwidth}


~\\
\textbf{\textsc{Course Objectives and Learning Outcomes}}
\begin{adjustwidth}{1cm}{0cm}
  After successfully completing this course, you will be able to:
  \begin{itemize}
    \item Design and calculate your own WAR metric using R and publicly available MLB data.
    \item Separate signal from noise in baseball performance data, and build a player projection model.
    \item Describe applications of analytics for supporting in-game strategy and player development.
    \item Identify an interesting research problem in baseball, solve it with data, and present your solution.
  \end{itemize}
\end{adjustwidth}

~\\
\textbf{\textsc{What You Need to Know}}
\begin{adjustwidth}{1cm}{0cm}
  \textbf{ThuRsday.} Thursday is R day. Please bring your laptop to class every Thursday because we will work through interactive R exercises to practice what we've been learning.

  \textbf{Exit Tickets and Quizzes.} We will use the final 15 minutes of almost every class to complete an exit ticket (anonymous) and a quiz (not anonymous). The purpose of the exit ticket is to collect feedback on the class. The purpose of the quiz is to practice retrieving information from memory, which is an effective way to learn things (according to pedagogical research). Quizzes are cumulative, not just based on the day's material.

  \textbf{Classmate Interaction.} In the middle of almost every class period, we will break out for $1$-on-$1$ discussions between classmates. The purpose is to give you a mental break and an opportunity to get to know your classmates. You are encouraged to find a different discussion partner every day.

  \textbf{Guest Speakers.} We plan to have $3$-$4$ guest speakers during the semester depending on availability.
\end{adjustwidth}

\newpage
\textbf{\textsc{How Your Grade Is Determined}}
\begin{adjustwidth}{1cm}{0cm}
	Sport analytics is a demanding and highly competitive field. This class aims to provide the tough but honest feedback that you will need to succeed in it.
  \begin{center}
    \begin{tabular}{rclcrclcrcl}
                      &               &     & & $[93\%,\,\infty)$ & $\rightarrow$ & A & & $[90\%, 93\%)$ & $\rightarrow$	& A--\\
      $[87\%, 90\%)$  & $\rightarrow$	& B+  & & $[83\%, 87\%)$  & $\rightarrow$ & B & & $[80\%, 83\%)$ & $\rightarrow$	& B--\\
      $[77\%, 80\%)$  & $\rightarrow$	& C+  & & $[73\%, 77\%)$  & $\rightarrow$ & C & & $[70\%, 73\%)$ & $\rightarrow$	& C--\\
      $[67\%, 70\%)$  & $\rightarrow$	& D+  & & $[63\%, 67\%)$  & $\rightarrow$ & D & & $[60\%, 63\%)$ & $\rightarrow$  & D--\\
                      &               &     & & $[ 0\%, 60\%)$  & $\rightarrow$ & F\\
    \end{tabular}
  \end{center}
  ~\\
	\textbf{Assignments.} ($40\%$) These are individual assignments. You are encouraged to help each other, but your submission must reflect work you have done yourself. Assignments will be graded anonymously, so please remove any personally identifying information from your submission.
  \begin{itemize}
    \item Assignment \#$1$ ($10\%$): Signal v. Noise in Batter Outcomes   \hfill  expected hours: $10$
    \item Assignment \#$2$ ($10\%$): Signal v. Noise in Pitcher Outcomes  \hfill  expected hours: $10$
    \item Assignment \#$3$ ($10\%$): Pitch Outcome Model                  \hfill  expected hours: $10$
    \item Assignment \#$4$ ($10\%$): WAR Projections                      \hfill  expected hours: $10$
  \end{itemize}

  \textbf{Project.} ($40\%$) This is a project-based course, and the project you build in stages over the course of the semester will determine a large chunk of your grade. Projects may be completed individually or in pairs (no more than two students).
  \begin{itemize}
    \item Project \#$0$ ($0\%$): Registration                   \hfill expected hours: $1$
    \item Project \#$1$ ($5\%$): Proposal                       \hfill expected hours: $5$
    \item Project \#$2$ ($10\%$): Abstract                      \hfill expected hours: $10$
    \item Project \#$3$ ($5\%$): Presentation                   \hfill expected hours: $5$
    \item Project \#$4$ ($20\%$): Paper                         \hfill expected hours: $20$
  \end{itemize}

  \textbf{Quizzes.} ($10\%$) We have a total of 24 quizzes planned for the semester. These are intended as low-stakes opportunities to test your learning, not as high-stakes assessments.

  \textbf{Classmate Interaction.} ($5\%$) We will have a total of approximately 24 opportunities for ``discussion break'' interactions during the semester. Your grade for this component will be based on the number of interactions you have, where you can only count up to four interactions with each different classmate. This is to incentivize you to mix up your discussion partners.

  \textbf{Late Work.} You may submit late work, subject to a $1$-point deduction for each late day ($24$-hour period).

  \textbf{Absence Policy.} Please notify me in advance if you are going to miss a class period.
\end{adjustwidth}

\newpage
\textbf{\textsc{Rice Honor Code}}
\begin{adjustwidth}{1cm}{0cm}
  In this course, all students will be held to the standards of the Rice Honor Code, a code that you pledged to honor when you matriculated at this institution. If you are unfamiliar with the details of this code and how it is administered, you should consult the Honor System Handbook at \url{http://honor.rice.edu/honor-system-handbook/}. This handbook outlines the University's expectations for the integrity of your academic work, the procedures for resolving alleged violations of those expectations, and the rights and responsibilities of students and faculty members throughout the process.\\
  ~\\
  \textbf{AI Policy.} You are allowed (and even encouraged) to use AI as a tool for developing R code (attribution in this case is not necessary). Please do not use AI to generate any part of your writeups.
\end{adjustwidth}

~\\
\textbf{\textsc{Disability Resource Center}}
\begin{adjustwidth}{1cm}{0cm}
  If you have a documented disability or other condition that may affect academic performance you should: $1$) make sure this documentation is on file with the Disability Resource Center (Allen Center, Room $111$ / \href{mailto:adarice@rice.edu}{adarice@rice.edu} / x$5841$) to determine the accommodations you need; and $2$) talk with me to discuss your accommodation needs.
\end{adjustwidth}

~\\
\textbf{\textsc{Mental Health Policy}}
\begin{adjustwidth}{1cm}{0cm}
	The wellbeing and mental health of students is important; if you are having trouble completing your coursework, please reach out to the Wellbeing and Counseling Center. Rice University provides cost-free mental health services through the Wellbeing and Counseling Center to help you manage personal challenges that threaten your personal or academic well-being. If you believe you are experiencing unusual amounts of stress, sadness, or anxiety, the Student Wellbeing Office or the Rice Counseling Center may be able to assist you. The Wellbeing and Counseling Center is located in the Gibbs Wellness Center and can be reached at 713-348-3311 (available 24/7).
\end{adjustwidth}

~\\
\textbf{\textsc{Title IX Responsible Employee Notification}}
\begin{adjustwidth}{1cm}{0cm}
  At Rice University, unlawful discrimination in any form, including sexual misconduct, is prohibited under Rice Policy on Harassment and Sexual Harassment (Policy 830) and the Student Code of Conduct.

  Please be aware that all employees of Rice University are ``mandatory reporters'', which means that if you tell me about a situation involving sexual harassment, sexual assault, dating violence, domestic violence, or stalking, I must share that information with the Title IX Coordinator.

  Although I have to make that notification, you will control how your case will be handled, including whether or not you wish to pursue a formal complaint. Our goal is to make sure you are aware of the range of options available to you and have access to the resources you need.

  To report sexual harassment, please contact the Title IX Coordinator at titleix@rice.edu. To explore supportive measures and other resources that are available to you, please visit the Office if Interpersonal Misconduct Prevention and Support at safe.rice.edu.
\end{adjustwidth}

\newpage

\textbf{\textsc{Consent for IRB Protocol FY2024-47}}
\begin{adjustwidth}{1cm}{0cm}
  This semester your instructor is participating in a study looking at teaching and learning. The purpose of this study is to determine the impact of the Center for Teaching Excellence's Instructor Network to Support Teaching Innovation and eLevate Learning (INSTILL) program. As part of that evaluation, your instructor will share data from this course with the CTE; no personally identifiable student information will be saved. Your participation in this study will last for the duration of the current semester and will entail activities no different from the regular activities you would otherwise engage in as part of the course. If you would like more information, please do not hesitate to ask.  If you would like to opt out of the study, please talk to your instructor as soon as possible.
\end{adjustwidth}

~\\
\textbf{This syllabus is only a guide for the course and is subject to change with advance notice.}

\end{document}
